%! TEX root = pertemuan_organisasi_ruang.tex
% the main TeX file which is intended to compile, :VimtexReload after adjustment
   % this file should be included in the main file
% :h vimtex-tex-root
\documentclass[../spa_2.tex]{subfiles}
%ensure the class of docs

\begin{document}

\begin{frame}[mytitle=standard,c,mycolor=digiPH_uniblue,light]
	\frametitle{Zonasi}

	Sebelum membahas zonasi, kita harus mengetahui terlebih dahulu: \textbf{Organisasi ruang}. Organisasi Ruang berfungsi untuk memahami bagaimana terjadinya hubungan antar ruang.

	Dengan kebutuhan ruang standar yang sudah ada, perencanaan dan pengorganisasian ruang dapat dibuat dengan membuat studi :

	\begin{itemize}
		\item Zoning
		\item  Grouping
		\item  Sirkulasi
	\end{itemize}\end{frame}



\begin{frame}[mytitle=standard,c,mybg=bglight3,light]
	\frametitle{Definisi Zonasi}
	\txt{digiPH_ocean}{digiPH_white}{Zoning merupakan  pembagian area secara umum berdasarkan jenis aktivitas yang dilakukan penghuni (dalam hal ini rumah tinggal) }

	\txt{digiPH_ocean}{digiPH_white}{Misalnya area yang diperuntukkan aktivitas dengan privasi tinggi akan semakin jauh dari area untuk umum (publik)    }
\end{frame}

\begin{frame}[mytitle=standard,c,mycolor=digiPH_gray,light]
	\frametitle{Macam-macam Zonasi}

	\begin{itemize}
		\item Zona Publik
		\item Zona Semi Publik
		\item Zona Semi Privat
		\item Zona Privat
		\item Servis
	\end{itemize}
\end{frame}



\begin{frame}[mytitle=standard,t,mycolor=digiPH_uniblue,light]
	\frametitle{Macam-macam Zonasi}
	\framesubtitle{Zona Publik}
	Zona ini merupakan wadah dari aktivitas umum yang digunakan sebagai ruang interaksi antara anggota keluarga dengan tamu atau orang luar yang diundang masuk ke rumah kita.

\end{frame}

\begin{frame}[mytitle=standard,c,mycolor=digiPH_leaf,dark]

	\frametitle{Macam-macam Zonasi}
	Yang termasuk dalam zona publik adalah :
	\begin{itemize}
		\item Lobi
		\item Ruang tamu
	\end{itemize}

	\txt{digiPH_white}{digiPH_black}{Pada umumnya zona publik ini mempunyai akses ke pintu utama dimana pencapaiannya mudah baik dari luar maupun dari dalam.}
\end{frame}


\begin{frame}[mytitle=standard,c,mybg=bglight2,light]

	\frametitle{Macam-macam Zonasi}
	\framesubtitle{Zona Semi Publik (Semi Public Zone)}

	\txt{digiPH_red}{digiPH_white}{Merupakan tempat sebagai wadah interaksi dengan seluruh anggota keluarga (penghuni rumah) dan kadang dengan  orang luar terdekat (kerabat)}

	Yang termasuk zona semi publik adalah :
	\begin{itemize}
		\item Kolam Renang
		\item Gym
	\end{itemize}
\end{frame}

\begin{frame}[mytitle=standard,c,mycolor=digiPH_ubcblue,dark]

	\frametitle{Macam-macam Zonasi}
	\framesubtitle{Zona Pribadi (Private Zone)}

	Merupakan zona yang biasanya hanya diakses dan digunakan oleh anggota keluarga (penghuni rumah). Yang termasuk dalam zona pribadi antara lain adalah :

	\begin{itemize}
		\item Kamar
		\item Kantor
	\end{itemize}
\end{frame}


\begin{frame}[mytitle=standard,c,mycolor=digiPH_cream,light]

	\frametitle{Macam-macam Zonasi}
	\framesubtitle{Zona Pribadi (Private Zone)}
	Zona pribadi biasanya diletakkan pada titik terjauh dari bagian muka bangunan, atau menjauhi dari daerah bising, polusi atau yang kurang nyaman seperti pinggir jalan raya.

	Yang perlu diperhatikan dalam perencanaan zona pribadi adalah :

	\begin{itemize}
		\item didesain senyaman mungkin
		\item mempunyai bukaan langsung ke ruang terbuka, seperti taman belakang atau samping
	\end{itemize}
\end{frame}

\begin{frame}[mytitle=standard,c,mycolor=digiPH_uniblue,light]

	\frametitle{Macam-macam Zonasi}
	\framesubtitle{Zona Semi Privat (Semi Private Zone)}
	Zona ini tingkat privasinya masih tinggi namun orang lain/tamu dengan kepentingan tertentu dapat masuk ke area ini. Yang termasuk zona semi privat adalah :

	\begin{itemize}
		\item Ruang hobi
		\item Ruang musik
		\item Home theatre,
		\item Perpustakaan, dll
	\end{itemize}\end{frame}

\begin{frame}[mytitle=standard,c,mybg=bglight2,light]

	\frametitle{Macam-macam Zonasi}
	\framesubtitle{Zona Servis (Service Zone)}

	Zona servis ini berfungsi melayani seluruh zona yang ada di dalam rumah, biasanya berbentuk garis lurus memanjang dari depan ke belakang rumah  atau memutar di sekitar rumah kita.

	Yang termasuk dalam zona servis adalah :
	\begin{itemize}
		\item Kamar mandi
		\item dapur
		\item ruang tidur pembantu
		\item gudang
		\item ruang cuci, jemur, setrika
	\end{itemize}

\end{frame}



\begin{frame}[mytitle=standard,c,mycolor=digiPH_leaf,dark]

	\begin{center}
		Hal-hal yang perlu diperhatikan dalam perencanaan zona servis adalah :
	\end{center}

	\begin{itemize}
		\item Persyaratan penanggulanan pemadam kebakaran (agar mudah menjangkau seluruh sudut rumah)
		\item Sumber penyediaan air bersih
		\item Saluran pembuangan air kotor
	\end{itemize}

	Sebaiknya berada dalam 1 garis dalam zona ini sehingga ketika  terjadi kebocoran akan mudah ditemukan dan tidak banyak lantai yang harus dibongkar.
\end{frame}

\begin{frame}[mytitle=standard,c,mybg=bglight2,light]
	\frametitle{Diagram Buble}

	Merupakan skema keterkaitan antar ruangan, lanjutan dari organisasi hubungan antar ruang tapi sudah lebih mendetail untuk pencapaian dan faktor pendukungnya.
\end{frame}


\begin{frame}[mytitle=standard,c,mybg=bglight2,light]
	\frametitle{Diagram Buble}
	\begin{figure}
		\begin{center}
			\includegraphics[width=.82\textwidth]{diagram_buble}
			% filename may not contain space but _
		\end{center}
		\caption{Diagram Bubble}
		\label{fig:bubble}
	\end{figure}
\end{frame}


\section{Diagram buble}%
\label{sec:Diagram buble}

\begin{frame}[mytitle=standard,c,mycolor=digiPH_uniblue,light]
	\frametitle{Diagram Buble}
	\begin{center}
		Setiap bulatan merupakan perwakilan dari tiap ruangan yang ada dalam perencanaan, besaran bulatan  menentukan prioritas kegiatan (menjadi pusat aktivitas) dan garis-garis yang menghubungkan merupakan intensitas kebutuhan dari penghuni
	\end{center}
\end{frame}

\begin{frame}[mytitle=center,c,mybg=bglight2,light]
	\frametitle{Contoh Kasus (asrama)}
	\begin{center}
		Tugas
	\end{center}

	\begin{itemize}
		\item Diagram Matriks
		\item Diagram Zonasi
	\end{itemize}
\end{frame}
\bglightimg{organisasi_12}

\bglightimg{organisasi_2}
\bglightimg{organisasi_3}
\bglightimg{organisasi_4}
\bglightimg{organisasi_5}
\bglightimg{organisasi_6}
\bglightimg{organisasi_7}
\bglightimg{organisasi_8}
\bglightimg{organisasi_9}
\bglightimg{organisasi_10}




\bibliographystyle{apalike}
\bibliography{../filename.bib} % required for citecompletion, not work in mainfile
\end{document}

